\section{Energy completion of a ramified current family}
\label{rc:sec}

A denominator which suffices in each finite energy range need not
suffice for an infinite state.  For a parameter \(\varpi\) and a
positive integer \(r\), the sequence \(\varpi^{-rE}\), indexed by
\(E\geq1\), has a Laurent polynomial
in every coordinate, but no common bound on its negative powers.
This elementary distinction changes the generic cokernel of a
current comparison.  It also separates that comparison from a
relative complex which admits one uniform contracting denominator.

Throughout this section, \(k\) is a characteristic-zero field and
\[
 R=k[\varpi],\qquad R_\varpi=k[\varpi,\varpi^{-1}],
 \qquad K=k(\varpi).
\]
The parameter \(\varpi\) is independent of the normal
coordinates of Section~\ref{ns:sec} and of the field coordinates
used below.  We introduce energy products explicitly here.
The algebraic sums in Sections~\ref{wo:sec} and~\ref{ns:sec}
retain their original meaning.

\subsection{The algebraic current map}

Fix integers \(a>b\geq1\), put \(r=a-b\), and choose \(d\geq1\)
current pairs.  Let \(V_R^A\) be the free vacuum current vertex
algebra on even fields \(U_i,V_i\), with generating products
\begin{equation}\label{rc:eq:source-ope}
 U_i(x)V_j(y)\sim\frac{\varpi^a\delta_{ij}}{(x-y)^2}.
\end{equation}
All other generating singular products vanish.  Let \(V_R^B\)
have pairs \(u_i,v_i\), with first level \(\varpi^b\) and remaining
levels \(\varpi^a\).  Coefficients in \(R\) are translation-constant.

These algebras have a direct vacuum description, valid even where
the levels vanish.  For a pair of level \(q\in R\), the modes obey
\[
 [U_m,V_n]=m q\,\delta_{m+n,0},\qquad
 [U_m,U_n]=[V_m,V_n]=0.
\]
All nonnegative modes kill the vacuum.  Negative modes act freely
by multiplication on the polynomial module in their symbols.
For \(m>0\), the remaining actions are
\begin{equation}\label{rc:eq:positive-modes}
 U_m=mq\,\partial_{V_{-m}},\qquad
 V_m=mq\,\partial_{U_{-m}}.
\end{equation}
The translation satisfies
\(T(U_{-m})=mU_{-m-1}\) and \(T(V_{-m})=mV_{-m-1}\).
The current fields are mutually local of order two.
Their normally ordered products and derivatives reconstruct the
fields of polynomial states: a creator \(U_{-m}\) contributes
\(\partial_x^{m-1}U(x)/(m-1)!\), and likewise for \(V\).
Each such field acts with ordinary lower truncation on polynomial
states.  This is the algebraic vacuum construction; no state
completion is used in it.

Define
\begin{equation}\label{rc:eq:current-map}
 j(U_1)=\varpi^r u_1,\qquad j(V_1)=v_1,\qquad
 j(U_i)=u_i,\quad j(V_i)=v_i\quad(i\geq2).
\end{equation}
The source pairing is preserved because
\(\varpi^r\varpi^b=\varpi^a\).  The mode relations and
normally ordered reconstruction therefore extend \(j\) to a map
preserving the vacuum, translation and every vertex product.

Index monomials in negative modes by finite exponent arrays
\(\lambda=(p_{i,m},q_{i,m})\).  Write \(e_\lambda,f_\lambda\)
for the corresponding source and target monomials, and set
\begin{equation}\label{rc:eq:energy}
 E(\lambda)=\sum_{i,m\geq1}m(p_{i,m}+q_{i,m}),\qquad
 \ell(\lambda)=\sum_{m\geq1}p_{1,m}.
\end{equation}
Thus \(E\) is energy and \(\ell\) counts active \(U_1\)-creators.
Every energy piece is finite free over \(R\), since only modes
\(m\leq E\) and finitely many exponent arrays can occur.  Moreover
\(0\leq\ell\leq E\); the monomial \(U_{1,-1}^E\mathbf1\)
attains \(\ell=E\).  This grading uses no conformal vector or
positive invariant form.

\begin{proposition}[The algebraic diagonal]\label{rc:prop:diagonal}
The map \(j\) is injective and satisfies
\begin{equation}\label{rc:eq:diagonal}
 j(e_\lambda)=\varpi^{r\ell(\lambda)}f_\lambda.
\end{equation}
Its cokernel as an \(R\)-module is
\begin{equation}\label{rc:eq:sum-cokernel}
 \bigoplus_{\ell(\lambda)>0}
       R/(\varpi^{r\ell(\lambda)})f_\lambda.
\end{equation}
After inverting \(\varpi\), it is an isomorphism of the algebraic
current vertex algebras.
\end{proposition}
\begin{proof}
Each occurrence of an active creator contributes the factor in
\eqref{rc:eq:current-map}; the other creators contribute one.
This proves \eqref{rc:eq:diagonal}.  The monomial bases are free,
and their diagonal entries are nonzero in the domain \(R\).
Coordinatewise quotient gives \eqref{rc:eq:sum-cokernel}.
Each vector in that direct sum has finite support, hence is
annihilated by some power of \(\varpi\).  Its localization is zero.
The inverse on localized generators sends \(u_1\) to
\(\varpi^{-r}U_1\) and leaves the other generators unchanged.
It preserves the generating pairings and therefore every
algebraic vertex product.
\end{proof}

\subsection{Two orders of completion and localization}

Give every coefficient module the discrete topology.  For
\(\bullet=A,B\), define
\begin{equation}\label{rc:eq:three-carriers}
 \widehat V_R^\bullet=\prod_{E\geq0}V_{R,E}^\bullet,\qquad
 L^\bullet=\widehat V_R^\bullet[1/\varpi],\qquad
 P^\bullet=\prod_{E\geq0}(V_{R,E}^\bullet[1/\varpi]).
\end{equation}
The first and third products have their product topologies.
The localization defining \(L^\bullet\) is algebraic.
It is not a completed localization.

\begin{lemma}[Common denominators]\label{rc:lem:denominators}
The natural map \(c^\bullet:L^\bullet\to P^\bullet\) is injective.
Its image consists precisely of the sequences whose Laurent
coordinates have a common bound on their negative powers of
\(\varpi\).
\end{lemma}
\begin{proof}
A localized product element is \(\varpi^{-N}v\) for one integer
\(N\geq0\) and one product vector \(v\).  This gives the asserted
bound.  Conversely, a common bound lets all coordinates be
multiplied by the same \(\varpi^N\), giving an element of
\(\widehat V_R^\bullet\).  If the image of
\(\varpi^{-N}v\) is zero, every coordinate of \(v\) is zero,
because the finite free energy modules are torsion-free.
\end{proof}

The diagonal gives a commutative diagram
\begin{equation}\label{rc:eq:comparison-square}
\begin{tikzcd}
 L^A \arrow[r,"j_L"] \arrow[d,"c^A"']&
 L^B \arrow[d,"c^B"]\\
 P^A \arrow[r,"j_P"']&P^B.
\end{tikzcd}
\end{equation}
The map \(j_P\) is a product-topological isomorphism.  Its inverse
divides the \(\lambda\)-coordinate by
\(\varpi^{r\ell(\lambda)}\), so each output coordinate depends
only on the corresponding input coordinate.

To compute the top arrow, put
\begin{equation}\label{rc:eq:product-cokernel}
 \mathcal T=
 \prod_{E\geq0}
 \bigoplus_{\substack{E(\lambda)=E\\\ell(\lambda)>0}}
       R/(\varpi^{r\ell(\lambda)})f_\lambda,\qquad
 \mathcal Q=\mathcal T[1/\varpi].
\end{equation}
The inner sums are finite.  These symbols denote module
cokernels; no quotient vertex algebra is implicit.

\begin{theorem}[A nonzero generic cokernel]\label{rc:thm:cokernel}
There are exact sequences
\begin{equation}\label{rc:eq:exact}
\begin{aligned}
 0&\longrightarrow\widehat V_R^A
   \xrightarrow{\widehat j}\widehat V_R^B
   \longrightarrow\mathcal T\longrightarrow0,\\
 0&\longrightarrow L^A\xrightarrow{j_L}L^B
   \longrightarrow\mathcal Q\longrightarrow0.
\end{aligned}
\end{equation}
The module \(\mathcal Q\) is a nonzero \(K\)-vector space.
In particular \(j_L\) is not surjective, although \(j_P\) is
an isomorphism.
\end{theorem}
\begin{proof}
Reduce every active coordinate modulo its displayed power of
\(\varpi\).  This is onto: polynomial representatives can be
chosen independently in each finite energy piece.  Its kernel
consists of the coordinatewise multiples of
\(\varpi^{r\ell(\lambda)}\).  Dividing each such coordinate
gives a source product vector, proving the first sequence.
Exactness of localization gives the second sequence.

Let \(w\in\widehat V_R^B\) have coefficient one on
\(f_E=u_{1,-1}^E\mathbf1\), for every \(E\geq1\), and zero
on all other monomials.  Its quotient class has coordinates
\begin{equation}\label{rc:eq:witness}
 q_E=1\bmod\varpi^{rE}.
\end{equation}
For each \(N\geq0\), choose \(E\) with \(rE>N\).
Then \(\varpi^Nq_E\ne0\), so \(q\) survives in \(\mathcal Q\).
Equivalently, the unique inverse of \(c^B(w)\) in \(P^A\)
has coordinates \(\varpi^{-rE}e_E\), with no common denominator.

If \(g\in k[\varpi]\) has \(g(0)\ne0\), multiplication by
\(g\) is an automorphism of every ring
\(R/(\varpi^{r\ell})\).  The coordinatewise inverses give
an automorphism of \(\mathcal T\), and then of \(\mathcal Q\).
Every nonzero polynomial is \(\varpi^Ng\) of this form.
Since \(\varpi\) is already invertible on \(\mathcal Q\),
every nonzero polynomial acts invertibly.  The universal
property of localization supplies the \(K\)-vector-space
structure.  The class \eqref{rc:eq:witness} proves it is nonzero.
No finite-dimensionality assertion is made.
\end{proof}

Unbounded torsion exponents alone would not prove this theorem:
the direct sum \eqref{rc:eq:sum-cokernel} has the same unbounded
exponents and zero localization.  The single product vector
\eqref{rc:eq:witness} is what defeats every exponent at once.
On the finite window \(E\leq M\), the cokernel is killed by
\(\varpi^{rM}\).  For any proposed uniform exponent \(N\),
the energy \(\lfloor N/r\rfloor+1\) disproves that bound in the
full product.

\begin{proposition}[Derived fibres away from the branch]
\label{rc:prop:fibres}
For every nonzero \(f\in R_\varpi\),
\begin{equation}\label{rc:eq:derived-fibre}
 \mathcal Q\otimes^{\mathbb L}_{R_\varpi}
                    R_\varpi/(f)=0.
\end{equation}
Thus derived base change of \(j_L\) is a quasi-isomorphism over
every such quotient.  Its generic base change still has the
nonzero cokernel \(\mathcal Q\).
\end{proposition}
\begin{proof}
For a nonunit \(f\), use the free resolution
\[
 [R_\varpi\xrightarrow{\ f\ }R_\varpi]
\]
in degrees \(-1,0\).  Tensoring gives
\([\mathcal Q\xrightarrow{f}\mathcal Q]\), whose differential
is invertible by Theorem~\ref{rc:thm:cokernel}.  An explicit
contraction sends the degree-zero term to the degree-minus-one
term by \(f^{-1}\).  For a unit \(f\) the quotient is zero.
Apply the derived tensor functor to the second exact sequence
in \eqref{rc:eq:exact}.  Finally
\(\mathcal Q\otimes_{R_\varpi}K=\mathcal Q\ne0\).
\end{proof}

Over \(k=\mathbb C\), this includes every nonzero closed fibre
\(\varpi=c\) and every finite thickening
\(f=(\varpi-c)^s\), \(s\geq1\).
The ordinary fibre has the expected product description:
\begin{equation}\label{rc:eq:ordinary-fibre}
 L^\bullet/(\varpi-c)L^\bullet
 \cong\prod_{E\geq0}
       \bigl(V_{R,E}^\bullet/(\varpi-c)V_{R,E}^\bullet\bigr).
\end{equation}
Indeed, coordinatewise evaluation is onto by constant
representatives.  Its kernel on the unlocalized product is
\((\varpi-c)\widehat V_R^\bullet\), by polynomial division
in each coordinate.  Since \(c\ne0\), localization changes
neither the quotient nor this assertion.  The specialized
diagonal has inverse \(c^{-r\ell(\lambda)}\) in every coordinate.
Thus the ordinary fibre statement and the derived assertion
both hold, without a finite-generation hypothesis on the family.

\begin{proposition}[The induced product topology]\label{rc:prop:dense}
Give \(L^\bullet\) the subspace topology from \(P^\bullet\)
through \(c^\bullet\).  Then \(j_L(L^A)\) is dense in \(L^B\).
The quotient topology on its nonzero algebraic cokernel
\(\mathcal Q\) is indiscrete, and its Hausdorff quotient is zero.
\end{proposition}
\begin{proof}
Every finite-energy-support target vector lifts under \(j_L\):
only finitely many extra powers of \(\varpi\) have to be
divided out.  Such vectors are dense in \(L^B\), since they
can match any prescribed finite set of coordinates.

More explicitly, let \(U\) be a basic neighbourhood of zero
requiring finitely many energy coordinates to vanish.
For each \(v\in L^B\), lift its finite truncation on those
coordinates.  The difference belongs to \(U\), so
\(U+j_L(L^A)=L^B\).  Therefore every nonempty open saturated
subset is all of \(L^B\), which proves the quotient-topology
claim.  Its Hausdorffization is zero.
\end{proof}

The topology in this proposition is part of its hypotheses.
The algebraic cokernel in Theorem~\ref{rc:thm:cokernel}
does not become a nonzero Hausdorff cokernel for this topology.
A different topology on algebraic localization would require
a separate definition and comparison.

\subsection{A relative complex with a uniform contraction}

Consider now the different polynomial and ghost complexes
\begin{equation}\label{rc:eq:relative-complexes}
\begin{aligned}
 A&=R[Q_1,\ldots,Q_d]\otimes\Lambda(\Theta_1,\ldots,\Theta_d),
 &\mathrm d_AQ_j&=\varpi^a\Theta_j,\\
 B&=R[q]\otimes\Lambda(\theta),
 &\mathrm d_Bq&=\varpi^b\theta.
\end{aligned}
\end{equation}
Polynomials have cohomological degree zero and ghosts degree
one.  All generators have auxiliary weight one; the
differentials preserve that weight.  For integers \(w_j\geq0\)
with \(w_1=0\), define a dg-algebra map
\begin{equation}\label{rc:eq:relative-map}
 \Phi(Q_j)=\varpi^{r+w_j}q,\qquad
 \Phi(\Theta_j)=\varpi^{w_j}\theta.
\end{equation}
It is a chain map because \(r+b=a\).
Its auxiliary weight \(n\) is distinct from the current energy
\(E\), and its target has one polynomial/ghost pair.

Use the cohomological cone
\begin{equation}\label{rc:eq:cone}
 C^i=B^i\oplus A^{i+1},\qquad
 \mathrm d_C(u,v)=(\mathrm d_Bu+\Phi v,-\mathrm d_Av).
\end{equation}
This is a complex.  No multiplication on its underlying
direct sum is presumed.

\begin{theorem}[A single contracting denominator]\label{rc:thm:cone}
There is a weight-preserving degree-minus-one operator \(H\)
on \(C\) such that
\begin{equation}\label{rc:eq:homotopy}
 \mathrm d_CH+H\mathrm d_C=\varpi^a\id,\qquad H^2=0.
\end{equation}
The same identities hold on
\(\widehat C=\prod_{n\geq0}C_n\).
Consequently \(\widehat C[1/\varpi]\) is contracted by
\(\varpi^{-a}H\).
\end{theorem}
\begin{proof}
For \(n>0\), let
\[
 \iota_A=\sum_jQ_j\iota_{\Theta_j},\qquad
 \iota_B=q\iota_\theta,\qquad
 h_A=\frac{\iota_A}{n},\qquad
 h_B=\frac{\varpi^r\iota_B}{n}.
\]
The polynomial Euler identity
\(\mathrm d_{\mathrm{pol}}\iota+\iota\mathrm d_{\mathrm{pol}}=n\)
gives
\[
 \mathrm d_Ah_A+h_A\mathrm d_A=\varpi^a\id,\qquad
 \mathrm d_Bh_B+h_B\mathrm d_B=\varpi^a\id.
\]
Division by \(n\) is valid in characteristic zero.
On generators, and then by the graded Leibniz rule,
\(\Phi\iota_A=\varpi^r\iota_B\Phi\).  Thus
\(\Phi h_A=h_B\Phi\).  Set
\[
 H(u,v)=(h_Bu,-h_Av).
\]
The off-diagonal term in \(\mathrm d_CH+H\mathrm d_C\)
is \(h_B\Phi-\Phi h_A=0\), and its diagonal terms are the
two Euler identities.  The contractions square to zero.

At weight zero the map is the identity of \(R\), and its
cone is \([R\xrightarrow{1}R]\) in degrees \(-1,0\).
Take \(H=\varpi^a\) times its elementary contraction.
All formulas preserve weight, so the identities extend
coordinatewise to the product.  The one denominator
\(\varpi^a\) suffices in every weight.
\end{proof}

In particular \(C\) is a dg module over
\(R[\epsilon]\), where \(|\epsilon|=-1\),
\(\epsilon^2=0\), and \(\mathrm d\epsilon=\varpi^a\):
let \(\epsilon\) act by \(H\).
This is a module over a derived thickening, not an asserted
algebra structure on the cone.  The resulting generic
acyclicity concerns \eqref{rc:eq:relative-complexes};
it does not annihilate the different current cokernel
\(\mathcal Q\).

\begin{proposition}[Auxiliary Hilbert bound for the cone]
\label{rc:prop:cone-norm}
Specialize \(\varpi=c\in\mathbb C\) with \(0<|c|\leq1\).
Give the polynomial/exterior spaces the auxiliary norm in which
\(Q^\alpha/\sqrt{\alpha!}\), respectively
\(q^j/\sqrt{j!}\), and the exterior monomials are orthonormal.
Use the orthogonal sum on the cone.  On positive weight \(n\),
\[
 \|c^{-a}H\|_{C_n}\leq |c|^{-a}n^{-1/2}.
\]
On the Hilbert sum of all weights, the contraction is bounded
by \(|c|^{-a}\) and preserves the domain of the closed
blockwise differential.  For \(d\geq2\), every bounded
contraction has norm at least \(|c|^{-a}\).
\end{proposition}
\begin{proof}
In the stated norm, \(\iota\) is the adjoint of
\(\mathrm d_{\mathrm{pol}}\), and their anticommutator is total
occupation \(n\).  Therefore
\(\|\iota/n\|\leq n^{-1/2}\).
The additional target factor has magnitude \(|c|^r\leq1\).
At weight zero the localized contraction has norm one.
Taking the orthogonal sum gives the bound.

Define the closed differential to have domain
\[
 \left\{(v_n)\in\bigoplus_{n\geq0}^{\ell^2}C_n:
                  \sum_n\|\mathrm d_Cv_n\|^2<\infty\right\}.
\]
It is the closed direct sum of its finite-dimensional blocks.
For the bounded blockwise homotopy \(S=c^{-a}H\),
the identity \(\mathrm d_CS v=v-S\mathrm d_Cv\) proves that
\(S\) maps this domain to itself.

For the lower bound, choose a unit weight-one ghost vector
\(v=\sum_jv_j\Theta_j\) in the kernel of
\((c^{w_1},\ldots,c^{w_d})\).  Such a vector exists if \(d\geq2\).
The cone vector \((0,v)\) is closed of degree zero.
Every degree-minus-one preimage has source linear-polynomial
part \(-c^{-a}\sum_jv_jQ_j\), forced by its source differential.
Orthogonal projection to this part has norm one.  Applying
any contraction to the closed vector therefore has norm at
least \(|c|^{-a}\).  Mixing other weights cannot change the
forced weight-one coefficient.
\end{proof}

This positive norm is a chosen estimate on a specified complex.
It does not assert positivity of an invariant current state.
Its uniform pole order will differ from the energy-dependent
inverse bound for the full current map.

\subsection{Fixed modes on the formal energy products}

A completed state need not define an ordinary vertex field.
Nevertheless, each specified mode has a finite contribution
to each output energy.  This supplies concrete operations on
the formal products.

For homogeneous algebraic states, the mode convention is
\begin{equation}\label{rc:eq:mode-energy}
 \mu_m(V_{R,p}\otimes_RV_{R,q})
       \subset V_{R,p+q-m-1},\qquad
 \mu_m(a,b)=a_{(m)}b.
\end{equation}
A negative output energy means zero.
To verify \eqref{rc:eq:mode-energy}, give each current state
energy one.  A current mode with index \(m\) changes energy
by \(-m\), as follows from multiplication and
\eqref{rc:eq:positive-modes}.  Translation raises energy by
one.  Normally ordered reconstruction then gives the formula
for every polynomial state.

For either current family, define the completed coefficient tensor
by finite energy truncations:
\begin{equation}\label{rc:eq:completed-tensor}
\begin{split}
 \widehat V_R\widehat\otimes_R\widehat V_R
 &:=
 \varprojlim_{M,N}
 \left(\bigoplus_{p\leq M}V_{R,p}\right)
 \otimes_R
 \left(\bigoplus_{q\leq N}V_{R,q}\right)\\
 &\cong\prod_{p,q\geq0}(V_{R,p}\otimes_RV_{R,q}).
\end{split}
\end{equation}
Every finite truncation has the discrete topology.
For \(P\), use the same construction over \(R_\varpi\),
replacing each \(V_{R,p}\) by \(V_{R,p}[1/\varpi]\).

\begin{proposition}[Continuous fixed-mode operations]
\label{rc:prop:modes}
For every integer \(m\), \(\mu_m\) extends to a continuous
linear map from \eqref{rc:eq:completed-tensor} to
\(\widehat V_R\), and to the analogous map for \(P\).
The induced bilinear operations are jointly continuous.
They restrict to \(L\), with its topology from
Proposition~\ref{rc:prop:dense}.
The maps \(\widehat j,j_P,j_L\) intertwine every such fixed
mode, translation and the vacuum.
\end{proposition}
\begin{proof}
At output energy \(E\), only the pairs
\[
 p,q\geq0,\qquad p+q=E+m+1
\]
can contribute.  There are finitely many, or none if the
right side is negative.  Each coefficient map is between
finite free modules.  Thus an output coordinate is a finite
sum of coefficient maps, and depends on finitely many input
coordinates.  This defines the continuous linear map on the
product tensor, without an infinite sum in the coefficient
ring.  The canonical map from two product vectors to their
pairwise tensors is jointly continuous, giving the bilinear
assertion.  The argument over \(R_\varpi\) is identical.

For \(L\), write the two inputs as \(\varpi^{-N}a\) and
\(\varpi^{-M}b\), with \(a,b\) in the unlocalized product.
Bilinearity gives
\[
 \mu_m(\varpi^{-N}a,\varpi^{-M}b)
     =\varpi^{-N-M}\mu_m(a,b)\in L.
\]
Thus the operation restricts, and continuity follows from
the subspace topology.  No additional tensor completion
of \(L\) is needed for this assertion.
Translation sends energy \(E\) to \(E+1\), so it extends
continuously as well.  The algebraic intertwining identities
hold in each of the finitely many contributions to any
output coordinate.  The vacuum identity is unchanged.
\end{proof}

These operations make the products into a system of
integer-indexed bilinear modes with translation and vacuum
at the stated scope.  No infinite sum of mode identities,
ordinary lower truncation on completed states, or convergence
of analytic fields is asserted.  The lower-truncation
failure will be explicit below.

There is already a precise obstruction to making the
cokernel a vertex quotient.  The proper image \(j_L(L^A)\)
contains the vacuum and is closed under its source mode
operations.  It is not an ideal for all target-state modes:
an ideal containing the vacuum would contain
\(v_{(-1)}\mathbf1=v\) for every \(v\in L^B\), including
the vector \(w\) of Theorem~\ref{rc:thm:cokernel}.
This argument uses the valid fixed-mode vacuum identity;
it does not require the completed target to be an ordinary
vertex algebra.

\subsection{A logarithmic parameter connection}

The diagonal current map is also compatible with a specified
connection in the parameter.  Let \(N_{U_i}\) count occurrences
of negative \(U_i\)-modes in a monomial, and similarly define
\(N_{u_i}\).  Let \(\delta=\varpi\partial_\varpi\) differentiate
only coefficients.  Define
\begin{equation}\label{rc:eq:connections}
 \mathcal D_A=\delta+a\sum_{i=1}^dN_{U_i},\qquad
 \mathcal D_B=\delta+bN_{u_1}
                      +a\sum_{i=2}^dN_{u_i}.
\end{equation}
Each satisfies
\(\mathcal D(fv)=\delta(f)v+f\mathcal D(v)\) for \(f\in R\).
This is a logarithmic parameter connection with respect to
the displayed regular monomial lattice.

\begin{theorem}[Horizontality of all fixed modes]
\label{rc:thm:connection}
On the algebraic current families,
\begin{equation}\label{rc:eq:connection-mode}
 \mathcal D(a_{(m)}b)
       =(\mathcal Da)_{(m)}b+a_{(m)}\mathcal Db
       \quad(m\in\mathbb Z).
\end{equation}
The vacuum is horizontal, \(\mathcal D\) commutes with
translation, and
\begin{equation}\label{rc:eq:horizontal-map}
 \mathcal D_Bj=j\mathcal D_A.
\end{equation}
The identities extend to every fixed mode on
\(\widehat V_R,L,P\).
The module cokernels \(\mathcal T,\mathcal Q\) inherit the
corresponding parameter connection.  After inverting
\(\varpi\), the operators \(\varpi^{-1}\mathcal D\) are
ordinary connection operators.
\end{theorem}
\begin{proof}
For a single pair of level \(\varpi^h\), put
\(\mathcal D=\delta+hN_U\).
A negative \(U\)-mode increases the occupation count by one,
while a negative \(V\)-mode leaves it unchanged.  For \(m>0\),
formula \eqref{rc:eq:positive-modes} gives
\[
 [\delta,U_m]=hU_m,\qquad [N_U,U_m]=0,
\]
and
\[
 [\delta,V_m]=hV_m,\qquad [N_U,V_m]=-V_m.
\]
Zero modes vanish.  Hence for every integer \(m\),
\[
 [\mathcal D,U_m]=hU_m,\qquad
 [\mathcal D,V_m]=0.
\]
The same calculation applies independently to all pairs.
Translation changes mode indices but not their occupation
counts, and the vacuum has count zero.
Applying the commutator Leibniz rule to the normally ordered
reconstruction proves \eqref{rc:eq:connection-mode} for
all algebraic states.  The central term has the correct
derivative as well:
\(\delta(\varpi^h)=h\varpi^h\).

On the \(\lambda\)-coordinate, differentiating
\(\varpi^{r\ell(\lambda)}\) contributes \(r\ell(\lambda)\).
This compensates the change from \(a\ell(\lambda)\) to
\(b\ell(\lambda)\), since \(b+r=a\).
The transverse counts agree, proving
\eqref{rc:eq:horizontal-map}.

These operators preserve energy.  They act coordinatewise
on products, and extend to a localization by the usual
coefficient differentiation rule.  Formula
\eqref{rc:eq:connection-mode} then extends using the finite
contributions of Proposition~\ref{rc:prop:modes}.
The ideals \((\varpi^{r\ell})\) are stable under
\(\delta\), so the diagonal quotient inherits the connection.
Finally \(\varpi^{-1}\delta=\partial_\varpi\) after localization.
\end{proof}

The bare coefficient derivative would fail on the first
active creator:
\begin{equation}\label{rc:eq:bare-defect}
 \delta j(U_{1,-1}\mathbf1)-j\delta(U_{1,-1}\mathbf1)
      =r\varpi^r u_{1,-1}\mathbf1\ne0.
\end{equation}
The occupation residues in \eqref{rc:eq:connections}
cancel exactly this defect.  This is differentiation of the
level parameter, not of a normal coordinate or a field
position.  It supplies no identification with the
position-Weyl objects of Section~\ref{ns:sec}.

\subsection{Hilbert norms and rapid energy decay}

Fix \(k=\mathbb C\) and \(c\ne0\), and write \(\rho=|c|\).
On each specialized energy piece, declare the monomials
\[
 \frac{e_\lambda}
 {\sqrt{\prod_{i,m}p_{i,m}!\,q_{i,m}!}},
 \qquad
 \frac{f_\lambda}
 {\sqrt{\prod_{i,m}p_{i,m}!\,q_{i,m}!}}
\]
orthonormal.  These are auxiliary positive norms.  No
invariance, current adjoint relation, or physical positivity
is imposed.

\begin{proposition}[Exact energywise bounds]\label{rc:prop:norms}
For every \(E\geq0\),
\begin{equation}\label{rc:eq:energy-norms}
 \|j_c\|_E=\max\{1,\rho^{rE}\},\qquad
 \|j_c^{-1}\|_E=\max\{1,\rho^{-rE}\}.
\end{equation}
In particular \(j_c^{-1}\) is unbounded on the unweighted
Hilbert sum if \(0<\rho<1\); \(j_c\) is unbounded if
\(\rho>1\).
\end{proposition}
\begin{proof}
The normalized monomial bases retain the diagonal
\(c^{r\ell(\lambda)}\).  Since \(0\leq\ell\leq E\),
the asserted maxima bound every diagonal entry.
The pure active state attains \(\ell=E\), and
\(V_{1,-1}^E\mathbf1\) attains \(\ell=0\).
These prove both equalities, including \(E=0\).
Taking the supremum over energy proves the assertions
about the Hilbert sum.
\end{proof}

For \(0<\rho<1\), the full-current inverse already exceeds
the relative-cone bound \(\rho^{-a}\) in every energy
\(E>a/r\).  Thus Proposition~\ref{rc:prop:cone-norm}
cannot be used as an inverse estimate for the current map.

Define the rapid-energy space in either monomial frame by
\begin{equation}\label{rc:eq:rapid-space}
 \mathscr S=
 \left\{v=(v_E):
 p_s(v)^2=\sum_{E\geq0}e^{2sE}\|v_E\|^2<\infty
                       \text{ for every }s>0\right\}.
\end{equation}
The positive-integer seminorms are cofinal.  They give a
complete Fr\'echet space: a Cauchy sequence converges in
each weighted Hilbert space, and the limits agree under
their continuous inclusions.

\begin{theorem}[Rapid comparison at every nonzero parameter]
\label{rc:thm:rapid}
For every \(c\in\mathbb C^\times\), \(j_c\) is a continuous
isomorphism between the two rapid-energy spaces.  More
precisely,
\begin{equation}\label{rc:eq:rapid-shifts}
\begin{aligned}
 p_s(j_cv)&\leq
 p_{s+r\max\{\log|c|,0\}}(v),\\
 p_s(j_c^{-1}v)&\leq
 p_{s+r\max\{-\log|c|,0\}}(v).
\end{aligned}
\end{equation}
Both shifts are uniform on a fixed compact annulus in
\(\mathbb C^\times\).  No fixed inverse shift, even with
a finite multiplicative norm constant, works for all
sufficiently small nonzero parameters.
Every fixed generator mode is a continuous linear
operator on \(\mathscr S\).
\end{theorem}
\begin{proof}
Square the bounds \eqref{rc:eq:energy-norms}, multiply
by \(e^{2sE}\), and sum over \(E\).  This gives
\eqref{rc:eq:rapid-shifts}, hence continuous inverse maps.
On an annulus, both logarithmic shifts are bounded.

For the last uniformity assertion, prescribe a shift
\(M\geq0\) and choose \(c\) such that
\(r\log(1/|c|)>M\).
For a normalized pure active state of energy \(E\),
the ratio of \(p_s(j_c^{-1}v)\) to \(p_{s+M}(v)\) is
\[
 \exp\bigl((r\log(1/|c|)-M)E\bigr).
\]
It is unbounded as \(E\to\infty\).  Thus no finite norm
constant repairs that shift at this chosen parameter.

A fixed generator mode changes energy by its fixed
integer index.  On normalized monomials its coefficient
has magnitude at most \(C_m\sqrt{E+1}\), by multiplication
or differentiation in \eqref{rc:eq:positive-modes}.
For every \(\varepsilon>0\),
\(\sqrt{E+1}e^{-\varepsilon E}\) is bounded.
The fixed energy shift contributes only another constant
to \(p_s\), proving
\[
 p_s(J_mv)\leq C_{m,s,\varepsilon,c}\,
                    p_{s+\varepsilon}(v).
\]
This proves continuity for each generator mode individually.
\end{proof}

These estimates concern a fixed coefficient value and a
specified topology.  They do not remove the family module
\(\mathcal Q\), and they do not assert that products of
arbitrary rapid states remain rapid.

\begin{example}[Failure of lower truncation]\label{rc:ex:truncation}
In the target at \(c\ne0\), let
\[
 v_\infty=\sum_{m\geq1}e^{-m^2}v_{1,-m}\mathbf1,\qquad
 u=u_{1,-1}\mathbf1.
\]
The sum belongs to every seminorm in
\eqref{rc:eq:rapid-space}, since
\(\sum_{m\geq1}e^{-2m^2+2sm}<\infty\).
Nevertheless,
\begin{equation}\label{rc:eq:truncation}
 u_{(m)}v_\infty=mc^b e^{-m^2}\mathbf1\ne0
                       \qquad(m\geq1).
\end{equation}
Indeed, only the creator \(v_{1,-m}\) contributes to the
positive current mode \(u_m\).
Thus ordinary vertex-algebra lower truncation fails even
for one algebraic generator acting on a rapid state.
Over \(R\), the formal energy product has the same
counterexample with coefficient one instead of \(e^{-m^2}\);
the output is \(m\varpi^b\mathbf1\ne0\).
\end{example}

The finite-energy-support subspace retains the original
algebraic vertex algebra.  The formal products retain the
continuous fixed-mode operations of
Proposition~\ref{rc:prop:modes}, and the rapid space
retains the stated generator-mode bounds.
A different analytic field recipient requires its own
convergence and multiplication theorems.
None of these coefficient constructions identifies
geometric collision operations or constructs an ordered
associative chiral duality.
